\documentclass[11pt]{ctexart}
\usepackage[a4paper,margin=1in]{geometry}
\usepackage{longtable,booktabs}
\title{Geant4-Agent Casebank Regression Report (EN)}
\author{Automated Pipeline}
\date{2026-03-08\_121732}
\begin{document}
\maketitle
\section{Unified Pass/Fail Table}
\begin{longtable}{p{0.06\linewidth}p{0.08\linewidth}p{0.07\linewidth}p{0.08\linewidth}p{0.07\linewidth}p{0.07\linewidth}p{0.07\linewidth}p{0.07\linewidth}p{0.09\linewidth}p{0.06\linewidth}p{0.21\linewidth}}
\toprule
ID & Domain & Style & Completeness & ExpInit & ActInit & ExpFinal & ActFinal & Missing(I->F) & Result & Fail Reason \\
\midrule
MTM01 & medical & fuzzy & missing & FAIL & FAIL & PASS & PASS & 12->0 & PASS & - \\
MTM02 & aerospace & precise & missing & FAIL & FAIL & PASS & PASS & 3->0 & PASS & - \\
MTM03 & industrial\_ndt & precise & missing & FAIL & FAIL & PASS & PASS & 12->0 & PASS & - \\
MTM04 & physics & precise & missing & FAIL & FAIL & PASS & PASS & 8->0 & PASS & - \\
MTM05 & nuclear\_engineering & fuzzy & missing & FAIL & FAIL & PASS & PASS & 9->0 & PASS & - \\
MTM06 & security\_screening & fuzzy & missing & FAIL & FAIL & PASS & PASS & 11->0 & PASS & - \\
MTM07 & semiconductor & precise & missing & FAIL & FAIL & PASS & PASS & 3->0 & PASS & - \\
MTM08 & education & fuzzy & missing & FAIL & FAIL & PASS & PASS & 11->0 & PASS & - \\
MTM09 & environmental\_radiation & fuzzy & missing & FAIL & FAIL & PASS & PASS & 3->0 & PASS & - \\
MTM10 & high\_energy\_physics & precise & missing & FAIL & FAIL & PASS & PASS & 10->0 & PASS & - \\
MTM11 & materials\_update & precise & full & PASS & PASS & PASS & PASS & 0->0 & PASS & - \\
MTM12 & overwrite\_reject & precise & full & PASS & PASS & PASS & PASS & 0->0 & PASS & - \\
MTM13 & medical\_detector & fuzzy & missing & FAIL & FAIL & PASS & FAIL & 13->1 & FAIL & final\_complete\_mismatch \\
MTM14 & industrial\_panel & precise & missing & FAIL & FAIL & PASS & PASS & 15->0 & FAIL & expected\_final\_mismatch:materials.selected\_materials.0 \\
MTM15 & shielding\_stack & precise & missing & FAIL & FAIL & PASS & PASS & 14->0 & PASS & - \\
MTM16 & nested\_target & precise & missing & FAIL & FAIL & PASS & PASS & 13->0 & PASS & - \\
MTM17 & coaxial\_shield & fuzzy & missing & FAIL & FAIL & PASS & PASS & 14->0 & PASS & - \\
MTM18 & boolean\_cavity & precise & missing & FAIL & FAIL & PASS & FAIL & 17->1 & FAIL & final\_complete\_mismatch, expected\_final\_mismatch:materials.selected\_materials.0 \\
MTM19 & ring\_complete & precise & full & PASS & FAIL & PASS & FAIL & 1->1 & FAIL & initial\_complete\_mismatch, final\_complete\_mismatch \\
MTM20 & grid\_complete & precise & full & PASS & FAIL & PASS & FAIL & 1->1 & FAIL & initial\_complete\_mismatch, final\_complete\_mismatch, expected\_final\_mismatch:geometry.structure, expected\_final\_mismatch:materials.selected\_materials.0 \\
MTM21 & layered\_detector & fuzzy & missing & FAIL & FAIL & PASS & PASS & 14->0 & PASS & - \\
MTM22 & embedded\_sensor & precise & full & PASS & PASS & PASS & PASS & 0->0 & PASS & - \\
MTM23 & layered\_shield\_complete & precise & full & PASS & PASS & PASS & PASS & 0->0 & PASS & - \\
MTM24 & boolean\_complete & precise & full & PASS & FAIL & PASS & FAIL & 1->1 & FAIL & initial\_complete\_mismatch, final\_complete\_mismatch, expected\_final\_mismatch:materials.selected\_materials.0, expected\_final\_mismatch:source.type \\
\bottomrule
\end{longtable}
\section{Baseline Comparison}
\begin{itemize}
\item pass\_rate=0.750
\item initial\_alignment\_rate=0.875
\item final\_alignment\_rate=0.792
\item avg\_turns\_per\_case=2.458
\item missing\_reduction\_rate=1.000
\item confirm\_apply\_success\_rate=1.000
\item rollback\_after\_confirm\_rate=0.000
\item internal\_leak\_turn\_rate=0.051
\item repeated\_assistant\_case\_rate=0.000
\item stale\_confirm\_case\_rate=0.000
\item delta\_pass\_rate=-0.200
\item regressions=none
\item improvements=none
\end{itemize}
\appendix
\section{Appendix: Full Dialogues and Naturalized Internal Chains}
\subsection{Case MTM01 (medical)}
\textbf{Dialogue:}
\begin{itemize}
\item User: 我先要一个用于医学成像的初步辐照场景，细节还没定，你先告诉我还缺什么。
\item Assistant: 目前我们已经确定了使用盒体作为几何结构，接下来还需要补充一些具体参数：盒体在x、y、z三个方向的尺寸，以及为体积指定材料，并选择源的类型。
\item User: 采用 300 mm x 300 mm x 300 mm 的水箱，材料 ?；源用束流 gamma，能量 6 MeV，位置 (0,0,-400) mm，方向 +z；物理列表 QBBC；输出 root。
\item Assistant: 我注意到这可能会覆盖已经确认的设置。具体来说：几何结构类型将从“未知”改为“单盒体”，同时盒体 x 方向尺寸将被设置为 300.0。

如果你想应用这些修改，请直接回复“确认”；如果想保留原来的值，请回复“保持原值”。
\item User: 确认
\item Assistant: 配置已完成！目前已经确认的设置有：几何结构类型、盒体在 x、y、z 三个方向的尺寸。
\end{itemize}
\textbf{Internal Chain (Natural Language):}
\begin{itemize}
\item Turn 1: action=summarize\_progress, phase=geometry, complete=False.
\item Missing fields: geometry.params.module\_x, geometry.params.module\_y, geometry.params.module\_z, materials.volume\_material\_map, source.type, source.particle, source.energy, source.position, source.direction, physics.physics\_list, output.format, output.path.
\item LLM used=True; fallback=none; backend=llm\_semantic\_frame+runtime\_semantic.
\item Turn 2: action=confirm\_overwrite, phase=geometry, complete=False.
\item Missing fields: geometry.params.module\_x, geometry.params.module\_y, geometry.params.module\_z.
\item LLM used=True; fallback=none; backend=llm\_slot\_frame+runtime\_semantic.
\item Turn 3: action=finalize, phase=materials, complete=True.
\item Missing fields: none.
\item LLM used=True; fallback=none; backend=llm\_slot\_frame+runtime\_semantic+confirmed\_pending\_overwrite.
\end{itemize}
\subsection{Case MTM02 (aerospace)}
\textbf{Dialogue:}
\begin{itemize}
\item User: I need a first-pass proton irradiation setup for a silicon device; keep it simple for now.
\item Assistant: Here's a summary of your setup so far and what's still needed.

**Completed:**
*   **Geometry:** A box structure with its X and Y dimensions defined.
*   **Materials:** Materials have been selected and assigned to your geometry volume.
*   **Source:** A proton source with its energy is configured.

**Still Required:**
*   **Physics:** You need to select a physics list.
*   **Output:** You need to specify the output format and the file path to save the data.
\item User: Use a silicon box 50 mm x 50 mm x 5 mm, beam source proton 100 MeV from (0,0,-100) mm along +z, physics FTFP\_BERT, output json.
\item Assistant: You've requested a change that would overwrite an already confirmed value.

Pending changes:
*   Box size along X: 1000.0 mm → 50.0 mm
*   Box size along Y: 1000.0 mm → 50.0 mm

Please reply **'confirm'** to apply these changes, or **'keep existing'** to keep the current values.
\item User: confirm
\item Assistant: Your configuration is complete. The following items are confirmed: box size along Z, box size along Y, box size along X, and the output path.
\end{itemize}
\textbf{Internal Chain (Natural Language):}
\begin{itemize}
\item Turn 1: action=summarize\_progress, phase=materials, complete=False.
\item Missing fields: physics.physics\_list, output.format, output.path.
\item LLM used=True; fallback=none; backend=llm\_slot\_frame+runtime\_semantic.
\item Turn 2: action=confirm\_overwrite, phase=source, complete=False.
\item Missing fields: none.
\item LLM used=True; fallback=none; backend=llm\_slot\_frame+runtime\_semantic.
\item Turn 3: action=finalize, phase=physics, complete=True.
\item Missing fields: none.
\item LLM used=True; fallback=none; backend=llm\_semantic\_frame+runtime\_semantic+confirmed\_pending\_overwrite.
\end{itemize}
\subsection{Case MTM03 (industrial\_ndt)}
\textbf{Dialogue:}
\begin{itemize}
\item User: 我要做一个 gamma 穿透铁靶的小测试，但我还没把完整参数说完。
\item Assistant: 目前我们已确认了材料（铁）和粒子类型（伽马射线）。接下来还需要你提供几何结构的具体类型和尺寸，以及体积与材料的绑定方式。请继续补充这些信息，我们就能一步步完成配置了。
\item User: 请设为半径 40 mm、半长 80 mm 的铁圆柱，材料 ?；点源 gamma，能量 2 MeV，位置 (0,0,-200) mm，方向 +z；物理列表 Shielding；输出 csv。
\item Assistant: 配置已完成。当前已确认的设置包括：体积与材料绑定、输出路径、输出格式和物理列表。
\end{itemize}
\textbf{Internal Chain (Natural Language):}
\begin{itemize}
\item Turn 1: action=summarize\_progress, phase=geometry, complete=False.
\item Missing fields: geometry.structure, geometry.params.module\_x, geometry.params.module\_y, geometry.params.module\_z, materials.volume\_material\_map, source.type, source.energy, source.position, source.direction, physics.physics\_list, output.format, output.path.
\item LLM used=True; fallback=none; backend=llm\_slot\_frame+runtime\_semantic.
\item Turn 2: action=finalize, phase=materials, complete=True.
\item Missing fields: none.
\item LLM used=True; fallback=none; backend=llm\_slot\_frame+runtime\_semantic.
\end{itemize}
\subsection{Case MTM04 (physics)}
\textbf{Dialogue:}
\begin{itemize}
\item User: Gamma point source 1 MeV at (0,0,-100) toward +z; please build the rest.
\item Assistant: Geometry: still needs geometry structure, box size along x. Materials: still needs volume-to-material binding. Source: updated source type, particle type, source energy. Physics: still needs physics list. Output: still needs output format, output path.
\item User: Set copper box 1 m x 1 m x 1 m, physics FTFP\_BERT, output hdf5.
\item Assistant: Your configuration is complete. The following items are confirmed: volume-to-material binding, output path, output format, and physics list.
\end{itemize}
\textbf{Internal Chain (Natural Language):}
\begin{itemize}
\item Turn 1: action=summarize\_progress, phase=geometry, complete=False.
\item Missing fields: geometry.structure, geometry.params.module\_x, geometry.params.module\_y, geometry.params.module\_z, materials.volume\_material\_map, physics.physics\_list, output.format, output.path.
\item LLM used=True; fallback=none; backend=llm\_slot\_frame+runtime\_semantic.
\item Turn 2: action=finalize, phase=materials, complete=True.
\item Missing fields: none.
\item LLM used=True; fallback=none; backend=llm\_slot\_frame+runtime\_semantic.
\end{itemize}
\subsection{Case MTM05 (nuclear\_engineering)}
\textbf{Dialogue:}
\begin{itemize}
\item User: 我想做一个中子屏蔽的概念验证，现在先给一个缺参开头。
\item Assistant: 目前我们已经确定了几个关键部分：几何结构类型、材料、粒子类型和物理列表都已同步。接下来还需要补充盒体的尺寸（x、y、z方向）、源的类型和能量，以及输出相关的设置。
\item User: 使用 200 mm x 200 mm x 200 mm 的铅块，材料 ?；点源 neutron，能量 14 MeV，位置 (0,0,-50) mm，方向 +z；物理列表 QGSP\_BERT\_HP；输出 xml。
\item Assistant: 我检测到这会覆盖已经确认的值。待修改内容：几何结构类型：unknown -> single\_box；已选材料：['G4\_Cu'] -> ['G4\_Pb']。如果要应用这些修改，请直接回复“确认”；如果要保留原值，请直接说“保持原值”。
\item User: 确认
\item Assistant: 配置已完成。当前已确认的设置包括：体积与材料绑定、盒体 z 方向尺寸、盒体 y 方向尺寸、盒体 x 方向尺寸。
\end{itemize}
\textbf{Internal Chain (Natural Language):}
\begin{itemize}
\item Turn 1: action=summarize\_progress, phase=geometry, complete=False.
\item Missing fields: geometry.params.module\_x, geometry.params.module\_y, geometry.params.module\_z, source.type, source.energy, source.position, source.direction, output.format, output.path.
\item LLM used=True; fallback=none; backend=llm\_slot\_frame+runtime\_semantic.
\item Turn 2: action=confirm\_overwrite, phase=geometry, complete=False.
\item Missing fields: geometry.params.module\_x, geometry.params.module\_y, geometry.params.module\_z.
\item LLM used=True; fallback=none; backend=llm\_slot\_frame+runtime\_semantic.
\item Turn 3: action=finalize, phase=materials, complete=True.
\item Missing fields: none.
\item LLM used=True; fallback=none; backend=llm\_semantic\_frame+runtime\_semantic+confirmed\_pending\_overwrite.
\end{itemize}
\subsection{Case MTM06 (security\_screening)}
\textbf{Dialogue:}
\begin{itemize}
\item User: Need a quick cargo-screening toy setup with forward photons, but I have not fixed the target shape, material, source energy, or output format yet.
\item Assistant: Geometry: updated geometry structure; still needs box size along x, box size along y. Materials: still needs volume-to-material binding. Source: updated particle type; still needs source type, source energy. Physics: still needs physics list. Output: still needs output format, output path.
\item User: Use steel cylinder target radius 60 mm half-length 120 mm, beam gamma 3 MeV from (0,0,-500) mm to +z, physics Shielding, output root.
\item Assistant: This would overwrite an already confirmed value. Pending change: geometry structure: unknown -> single\_tubs; box size along x: None -> 1000.0. Reply 'confirm' to apply it, or say 'keep existing' to keep the current value.
\item User: confirm
\item Assistant: Configuration complete. The following items are confirmed: volume-to-material binding, half length, radius, and box size along z.
\end{itemize}
\textbf{Internal Chain (Natural Language):}
\begin{itemize}
\item Turn 1: action=summarize\_progress, phase=geometry, complete=False.
\item Missing fields: geometry.params.module\_x, geometry.params.module\_y, geometry.params.module\_z, materials.volume\_material\_map, source.type, source.energy, source.position, source.direction, physics.physics\_list, output.format, output.path.
\item LLM used=True; fallback=none; backend=llm\_slot\_frame+runtime\_semantic.
\item Turn 2: action=confirm\_overwrite, phase=geometry, complete=False.
\item Missing fields: geometry.params.module\_x, geometry.params.module\_y, geometry.params.module\_z.
\item LLM used=True; fallback=none; backend=llm\_slot\_frame+runtime\_semantic.
\item Turn 3: action=finalize, phase=materials, complete=True.
\item Missing fields: none.
\item LLM used=True; fallback=none; backend=llm\_semantic\_frame+runtime\_semantic+confirmed\_pending\_overwrite.
\end{itemize}
\subsection{Case MTM07 (semiconductor)}
\textbf{Dialogue:}
\begin{itemize}
\item User: 我想搭一个硅薄片的中子照射测试，但先只给一个简述。
\item Assistant: 好的，我们已同步了以下配置：
*   **几何**：盒体的类型与X、Y方向尺寸。
*   **材料**：已选定材料并完成绑定。
*   **源**：粒子类型、能量和位置。

接下来，还需要您补充以下信息才能完成配置：
*   物理列表
*   输出格式
*   输出路径
\item User: 使用 20 mm x 20 mm x 1 mm 的硅片，材料 ?；各向同性 neutron，能量 2 MeV，位置 (0,0,0) mm；物理列表 QGSP\_BERT\_HP；输出 json。
\item Assistant: 我注意到新的尺寸设置会覆盖已确认的数值。具体变动如下：
- 盒体 x 方向尺寸：1000.0 → 20.0
- 盒体 y 方向尺寸：1000.0 → 20.0

请确认您的选择：直接回复“确认”以应用修改，或回复“保持原值”以保留现有设置。
\item User: 确认
\item Assistant: 配置已完成。当前已确认的参数包括：源位置、源能量、源类型以及盒体的 z 方向尺寸。
\end{itemize}
\textbf{Internal Chain (Natural Language):}
\begin{itemize}
\item Turn 1: action=summarize\_progress, phase=materials, complete=False.
\item Missing fields: physics.physics\_list, output.format, output.path.
\item LLM used=True; fallback=none; backend=llm\_slot\_frame+runtime\_semantic.
\item Turn 2: action=confirm\_overwrite, phase=source, complete=False.
\item Missing fields: none.
\item LLM used=True; fallback=none; backend=llm\_slot\_frame+runtime\_semantic.
\item Turn 3: action=finalize, phase=physics, complete=True.
\item Missing fields: none.
\item LLM used=True; fallback=none; backend=llm\_semantic\_frame+runtime\_semantic+confirmed\_pending\_overwrite.
\end{itemize}
\subsection{Case MTM08 (education)}
\textbf{Dialogue:}
\begin{itemize}
\item User: For classroom demo, I want gamma attenuation in copper; you can ask only essential missing parts.
\item Assistant: Here's a summary of your setup so far:

**Completed:**
*   **Materials:** Copper is selected.
*   **Source:** Particle type is set to gamma.
*   **Physics:** A physics list has been chosen.

**Still needed:**
*   **Geometry:** Structure type and box dimensions (X, Y, Z).
*   **Materials:** Assignment of the material to your geometry volume.
*   **Source:** Type, energy, and position.
*   **Output:** Format and file path.
\item User: Okay, copper box 100 mm x 100 mm x 100 mm, point gamma 1 MeV at (0,0,-200) mm, direction +z, physics FTFP\_BERT, output csv.
\item Assistant: This would change a setting you've already confirmed.

**Change:** Physics list from "Shielding" to "FTFP\_BERT".

Please reply **'confirm'** to apply this change, or **'keep existing'** to keep the current "Shielding" list.
\item User: confirm
\item Assistant: Your configuration is complete. The following items are confirmed: physics list, output path, output format, and source direction.
\end{itemize}
\textbf{Internal Chain (Natural Language):}
\begin{itemize}
\item Turn 1: action=summarize\_progress, phase=geometry, complete=False.
\item Missing fields: geometry.structure, geometry.params.module\_x, geometry.params.module\_y, geometry.params.module\_z, materials.volume\_material\_map, source.type, source.energy, source.position, source.direction, output.format, output.path.
\item LLM used=True; fallback=none; backend=llm\_slot\_frame+runtime\_semantic.
\item Turn 2: action=confirm\_overwrite, phase=materials, complete=False.
\item Missing fields: none.
\item LLM used=True; fallback=none; backend=llm\_slot\_frame+runtime\_semantic.
\item Turn 3: action=finalize, phase=source, complete=True.
\item Missing fields: none.
\item LLM used=True; fallback=none; backend=llm\_semantic\_frame+runtime\_semantic+confirmed\_pending\_overwrite.
\end{itemize}
\subsection{Case MTM09 (environmental\_radiation)}
\textbf{Dialogue:}
\begin{itemize}
\item User: 我想先搭一个空气体积中的低能 gamma 传播场景，细节稍后补。
\item Assistant: 好的，目前配置进度如下：

**已设置完成：**
*   **几何**：已确定使用盒体，并设置了X、Y方向的尺寸。
*   **材料**：已选定材料，并完成了与体积的绑定。
*   **源**：已设定粒子类型、能量和位置。

**仍需补充以下信息才能生成脚本：**
*   **物理**：请选择物理列表。
*   **输出**：请指定输出格式和文件保存路径。

请提供这些信息以继续。
\item User: 设为 500 mm x 500 mm x 500 mm 的空气盒，材料 ??；点源 gamma，能量 0.5 MeV，位置 (0,0,-300) mm，方向 +z；物理列表 QBBC；输出 root。
\item Assistant: 我注意到这会覆盖已确认的设置。具体修改如下：

- **盒体 x 方向尺寸**：1000.0 → 500.0
- **盒体 y 方向尺寸**：1000.0 → 500.0

请确认您的选择：
- 如果应用这些修改，请直接回复“确认”。
- 如果保留原值，请直接回复“保持原值”。
\item User: 确认
\item Assistant: 好的，配置已完成。当前已确认的设置包括：源位置、源能量、盒体在Z、Y、X三个方向的尺寸。
\end{itemize}
\textbf{Internal Chain (Natural Language):}
\begin{itemize}
\item Turn 1: action=summarize\_progress, phase=materials, complete=False.
\item Missing fields: physics.physics\_list, output.format, output.path.
\item LLM used=True; fallback=none; backend=llm\_slot\_frame+runtime\_semantic.
\item Turn 2: action=confirm\_overwrite, phase=source, complete=False.
\item Missing fields: none.
\item LLM used=True; fallback=none; backend=llm\_slot\_frame+runtime\_semantic.
\item Turn 3: action=finalize, phase=physics, complete=True.
\item Missing fields: none.
\item LLM used=True; fallback=none; backend=llm\_semantic\_frame+runtime\_semantic+confirmed\_pending\_overwrite.
\end{itemize}
\subsection{Case MTM10 (high\_energy\_physics)}
\textbf{Dialogue:}
\begin{itemize}
\item User: Set geometry as a sphere target; I have not decided radius or source details yet.
\item Assistant: Here's where we stand:

**Completed:**
*   **Geometry:** You've set up a sphere target and provided its dimensions (box sizes for the world volume).

**Still Needed:**
*   **Geometry:** The sphere's radius.
*   **Materials:** Which material to assign to the sphere.
*   **Source:** The source type, particle type, and its energy.
*   **Physics:** The physics list to use.
*   **Output:** The output format and file path.

What would you like to provide next?
\item User: Sphere target radius 80 mm, lead material, electron point source 5 MeV at (0,0,-200) mm toward +z, physics FTFP\_BERT, output hdf5.
\item Assistant: Your configuration is complete. All required settings, including the volume-to-material binding, output path, output format, and physics list, have been confirmed.
\end{itemize}
\textbf{Internal Chain (Natural Language):}
\begin{itemize}
\item Turn 1: action=summarize\_progress, phase=geometry, complete=False.
\item Missing fields: geometry.params.child\_rmax, materials.volume\_material\_map, source.type, source.particle, source.energy, source.position, source.direction, physics.physics\_list, output.format, output.path.
\item LLM used=True; fallback=none; backend=llm\_slot\_frame+runtime\_semantic.
\item Turn 2: action=finalize, phase=materials, complete=True.
\item Missing fields: none.
\item LLM used=True; fallback=none; backend=llm\_slot\_frame+runtime\_semantic.
\end{itemize}
\subsection{Case MTM11 (materials\_update)}
\textbf{Dialogue:}
\begin{itemize}
\item User: Create a complete setup: 100 mm x 100 mm x 100 mm copper box, point gamma 1 MeV at (0,0,-200) mm toward +z, physics QBBC, output json.
\item Assistant: Your configuration is complete. The following items have been confirmed: volume-to-material binding, output path, and output format.
\item User: Change the target material to aluminum and change the output format to root.
\item Assistant: You've requested changes that would overwrite already confirmed values:

*   **Selected materials:** Copper → Aluminum
*   **Output format:** JSON → ROOT

Please reply **'confirm'** to apply these changes, or **'keep existing'** to keep the current values.
\item User: confirm
\item Assistant: Your configuration is complete. The following items have been confirmed: selected materials, output format, output path, and volume-to-material binding.
\end{itemize}
\textbf{Internal Chain (Natural Language):}
\begin{itemize}
\item Turn 1: action=finalize, phase=materials, complete=True.
\item Missing fields: none.
\item LLM used=True; fallback=none; backend=llm\_slot\_frame+runtime\_semantic.
\item Turn 2: action=confirm\_overwrite, phase=source, complete=False.
\item Missing fields: none.
\item LLM used=True; fallback=none; backend=llm\_slot\_frame+runtime\_semantic.
\item Turn 3: action=finalize, phase=physics, complete=True.
\item Missing fields: none.
\item LLM used=True; fallback=none; backend=llm\_semantic\_frame+runtime\_semantic+confirmed\_pending\_overwrite.
\end{itemize}
\subsection{Case MTM12 (overwrite\_reject)}
\textbf{Dialogue:}
\begin{itemize}
\item User: 请创建一个完整配置：100 mm x 100 mm x 100 mm 的铜盒，材料 ?；点源 gamma，能量 1 MeV，位置 (0,0,-200) mm，方向 +z；物理列表 QBBC；输出 json。
\item Assistant: 配置已完成。当前已确认的设置包括：体积与材料绑定、输出路径和输出格式。
\item User: 把材料改成 ?，并把输出格式改成 root。
\item Assistant: 我注意到这次修改会覆盖已确认的设置。具体来说：输出格式将从 json 改为 root。

请确认你的选择：直接回复“确认”来应用修改，或回复“保持原值”以保留原来的设置。
\item User: 保持原值
\item Assistant: 好的，输出格式仍将保持为 json，没有更改。
\end{itemize}
\textbf{Internal Chain (Natural Language):}
\begin{itemize}
\item Turn 1: action=finalize, phase=materials, complete=True.
\item Missing fields: none.
\item LLM used=True; fallback=none; backend=llm\_slot\_frame+runtime\_semantic.
\item Turn 2: action=confirm\_overwrite, phase=source, complete=False.
\item Missing fields: geometry.ask.boolean.solid\_a\_size, geometry.ask.boolean.solid\_b\_size.
\item LLM used=True; fallback=none; backend=llm\_slot\_frame+runtime\_semantic.
\item Turn 3: action=reject\_overwrite, phase=physics, complete=True.
\item Missing fields: none.
\item LLM used=True; fallback=none; backend=llm\_semantic\_frame+runtime\_semantic+rejected\_pending\_overwrite.
\end{itemize}
\subsection{Case MTM13 (medical\_detector)}
\textbf{Dialogue:}
\begin{itemize}
\item User: I want a ring of detector modules around a small source region, but I have not fixed the ring radius or module size yet.
\item Assistant: Your geometry structure is now set for a ring of detector modules around a small source region.

Here’s what’s been completed and what is still needed next:

**Completed:**
*   **Geometry:** The structure (a ring of modules) is confirmed. The overall bounding box size is also set.

**Still Needed:**
*   **Geometry:** Ring module size, ring module count, and ring radius.
*   **Materials:** Which material to assign to each volume.
*   **Source:** Source type, particle type, and its energy and position.
*   **Physics:** Which physics list to use.
*   **Output:** The output format and file path.
\item User: Use a ring of 12 modules, each 8 mm x 10 mm x 2 mm, radius 40 mm, clearance 1 mm, material water; point gamma source 0.662 MeV at (0,0,0) mm; physics QBBC; output json.
\item Assistant: This would overwrite a geometry value you already confirmed.

**Pending Change:**
*   **Module Size:** From (7.83, 4.72, 6.86) mm to (8, 10, 2) mm.
*   **Ring Modules:** From 4 to 12.
*   **Ring Radius:** From 67.49 mm to 40 mm.
*   **Clearance:** From 0.19 mm to 1 mm.

Reply **'confirm'** to apply these changes, or **'keep existing'** to keep the current values.
\item User: confirm
\item Assistant: To confirm the source direction, what are the x, y, and z components (dx, dy, dz) for the beam direction?
\end{itemize}
\textbf{Internal Chain (Natural Language):}
\begin{itemize}
\item Turn 1: action=summarize\_progress, phase=materials, complete=False.
\item Missing fields: materials.volume\_material\_map, source.type, source.particle, source.energy, source.position, source.direction, physics.physics\_list, output.format, output.path, geometry.ask.ring.module\_size, geometry.ask.ring.count, geometry.ask.ring.radius, geometry.ask.ring.clearance.
\item LLM used=True; fallback=none; backend=llm\_semantic\_frame+runtime\_semantic.
\item Turn 2: action=confirm\_overwrite, phase=source, complete=False.
\item Missing fields: source.direction.
\item LLM used=True; fallback=none; backend=llm\_slot\_frame+runtime\_semantic.
\item Turn 3: action=ask\_clarification, phase=source, complete=False.
\item Missing fields: source.direction.
\item LLM used=True; fallback=none; backend=llm\_semantic\_frame+runtime\_semantic+confirmed\_pending\_overwrite.
\end{itemize}
\subsection{Case MTM14 (industrial\_panel)}
\textbf{Dialogue:}
\begin{itemize}
\item User: 我想做一个二维阵列探测板，但 pitch 和材料还没定。
\item Assistant: 目前我们已确定了使用二维阵列探测板作为几何结构类型，并完成了盒体尺寸的设置。

接下来还需要补充以下信息：
- **材料**：体积与材料绑定
- **源**：源类型、粒子类型、源能量、源位置
- **物理**：物理列表
- **输出**：输出格式、输出路径

请继续提供这些信息，以便完成配置。
\item User: ?? 4 x 3 ?????????? 10 mm x 10 mm x 2 mm?pitch\_x 12 mm?pitch\_y 14 mm??? 1 mm??? ???? gamma 0.511 MeV??? (0,0,-80) mm??? +z????? QBBC??? root?
\item Assistant: 我检测到这会覆盖已经确认的值。待修改内容：几何结构类型：grid -> single\_box；盒体 x 方向尺寸：None -> 10.0。如果要应用这些修改，请直接回复“确认”；如果要保留原值，请直接说“保持原值”。
\item User: 确认
\item Assistant: 配置已完成。当前已确认的参数包括：y 方向 pitch、x 方向 pitch、盒体 z 方向尺寸、盒体 y 方向尺寸。
\end{itemize}
\textbf{Internal Chain (Natural Language):}
\begin{itemize}
\item Turn 1: action=summarize\_progress, phase=materials, complete=False.
\item Missing fields: materials.volume\_material\_map, source.type, source.particle, source.energy, source.position, source.direction, physics.physics\_list, output.format, output.path, geometry.ask.grid.module\_size, geometry.ask.grid.count\_x, geometry.ask.grid.count\_y, geometry.ask.grid.pitch\_x, geometry.ask.grid.pitch\_y, geometry.ask.grid.clearance.
\item LLM used=True; fallback=none; backend=llm\_semantic\_frame+runtime\_semantic.
\item Turn 2: action=confirm\_overwrite, phase=source, complete=False.
\item Missing fields: none.
\item LLM used=True; fallback=none; backend=llm\_slot\_frame+runtime\_semantic.
\item Turn 3: action=finalize, phase=physics, complete=True.
\item Missing fields: none.
\item LLM used=True; fallback=none; backend=llm\_semantic\_frame+runtime\_semantic+confirmed\_pending\_overwrite.
\end{itemize}
\subsection{Case MTM15 (shielding\_stack)}
\textbf{Dialogue:}
\begin{itemize}
\item User: Build a stacked absorber sandwich along z; I will provide the layer sizes next.
\item Assistant: Your geometry structure is now set to a stacked absorber sandwich along the z-axis. Here's what's still needed:

**Geometry:** Stack layer footprint and thicknesses.
**Materials:** Which volumes use which materials.
**Source:** Type, particle, and energy.
**Physics:** Physics list selection.
**Output:** Format and file path.
\item User: Use stacked layers with footprint 20 mm x 20 mm, thicknesses 2 mm, 4 mm, 6 mm, layer clearance 0.5 mm, outer box 30 mm x 30 mm x 20 mm, container clearance 1 mm, material aluminum; beam gamma 1.25 MeV from (0,0,-120) mm along +z; physics Shielding; output csv.
\item Assistant: You are about to overwrite a previously confirmed geometry configuration.

**Pending Change:**
- **Stack footprint:** 16.33 mm x 10.28 mm → 20.0 mm x 20.0 mm
- **Layer thicknesses:** [6.68, 1.19, 5.59] mm → [2.0, 4.0, 6.0] mm
- **Layer clearance:** 0.73 mm → 0.5 mm
- **Outer box size:** 18.07 mm x 45.52 mm x 16.66 mm → 30.0 mm x 30.0 mm x 20.0 mm
- **Container clearance:** 0.87 mm → 1.0 mm

Reply **'confirm'** to apply this change, or **'keep existing'** to keep the current values.
\item User: confirm
\item Assistant: Configuration complete. The following geometry items are confirmed: stack container clearance, stack container size, stack layer clearance, and stack layer thicknesses.
\end{itemize}
\textbf{Internal Chain (Natural Language):}
\begin{itemize}
\item Turn 1: action=summarize\_progress, phase=materials, complete=False.
\item Missing fields: materials.volume\_material\_map, source.type, source.particle, source.energy, source.position, source.direction, physics.physics\_list, output.format, output.path, geometry.ask.stack.footprint, geometry.ask.stack.thicknesses, geometry.ask.stack.layer\_clearance, geometry.ask.stack.parent\_size, geometry.ask.stack.nest\_clearance.
\item LLM used=True; fallback=none; backend=llm\_semantic\_frame+runtime\_semantic.
\item Turn 2: action=confirm\_overwrite, phase=source, complete=False.
\item Missing fields: none.
\item LLM used=True; fallback=none; backend=llm\_slot\_frame+runtime\_semantic.
\item Turn 3: action=finalize, phase=physics, complete=True.
\item Missing fields: none.
\item LLM used=True; fallback=none; backend=llm\_semantic\_frame+runtime\_semantic+confirmed\_pending\_overwrite.
\end{itemize}
\subsection{Case MTM16 (nested\_target)}
\textbf{Dialogue:}
\begin{itemize}
\item User: 我想在一个盒子里嵌一个圆柱靶，先别急着补全。
\item Assistant: 好的，我们已确认了您要构建的“盒中嵌圆柱”几何结构。

目前，几何部分已明确结构类型，并确认了外层盒体的尺寸。接下来，我们还需要您提供：
1.  内嵌圆柱的尺寸（半径和半长）。
2.  为盒体和圆柱分别指定材料。
3.  设置粒子源的类型、种类和能量等信息。

请随时告诉我这些内容，我们继续配置。
\item User: 外盒 120 mm x 120 mm x 120 mm，内嵌铜圆柱半径 20 mm、半长 40 mm，嵌套间隙 2 mm；材料 ?；束流 gamma 1 MeV，位置 (0,0,-150) mm，方向 +z；物理列表 QBBC；输出 json。
\item Assistant: 我注意到这将会覆盖已确认的几何参数。具体来说，外层盒子的尺寸、内嵌圆柱的半径与半长，以及嵌套间隙都将被更新为新值。

如果您确认应用这些修改，请直接回复“确认”；如果您希望保留原来的数值，请回复“保持原值”。
\item User: 确认
\item Assistant: 好的，配置已完成。当前已确认的参数包括：嵌套间隙、内嵌圆柱半长、内嵌圆柱半径以及外层容器的尺寸。
\end{itemize}
\textbf{Internal Chain (Natural Language):}
\begin{itemize}
\item Turn 1: action=summarize\_progress, phase=materials, complete=False.
\item Missing fields: materials.volume\_material\_map, source.type, source.particle, source.energy, source.position, source.direction, physics.physics\_list, output.format, output.path, geometry.ask.nest.parent\_size, geometry.ask.nest.child\_radius, geometry.ask.nest.child\_half\_length, geometry.ask.nest.clearance.
\item LLM used=True; fallback=none; backend=llm\_slot\_frame+runtime\_semantic.
\item Turn 2: action=confirm\_overwrite, phase=source, complete=False.
\item Missing fields: none.
\item LLM used=True; fallback=none; backend=llm\_slot\_frame+runtime\_semantic.
\item Turn 3: action=finalize, phase=physics, complete=True.
\item Missing fields: none.
\item LLM used=True; fallback=none; backend=llm\_semantic\_frame+runtime\_semantic+confirmed\_pending\_overwrite.
\end{itemize}
\subsection{Case MTM17 (coaxial\_shield)}
\textbf{Dialogue:}
\begin{itemize}
\item User: I need a shell-style shielding toy model, but the shell thicknesses are still open.
\item Assistant: Here's what we have so far:
*   **Geometry:** The structure is set to a shell. We still need the inner radius and the thicknesses for each shell layer.
*   **Materials:** We need to assign materials to each volume.
*   **Source:** We need the source type, particle type, energy, and position.
*   **Physics:** The physics list is configured.
*   **Output:** We need to choose an output format and file path.
\item User: Use a shell with inner radius 15 mm, thicknesses 5 mm and 8 mm, half length 40 mm, nested core radius 10 mm, nested core half length 30 mm, clearance 1 mm, material lead; isotropic neutron source 2 MeV at (0,0,0) mm; physics Shielding; output hdf5.
\item Assistant: You're about to overwrite a geometry that was already set.

**Current geometry:**
*   **Shell:** Inner radius 7.15 mm, thicknesses [1.18, 3.43, 0.86] mm, half-length 22.36 mm.
*   **Nested Core:** Radius 11.60 mm, half-length 2.68 mm.
*   **Clearance:** 1.01 mm.

**New geometry:**
*   **Shell:** Inner radius 15.0 mm, thicknesses [5.0, 8.0, 0.83] mm, half-length 40.0 mm.
*   **Nested Core:** Radius 10.0 mm, half-length 30.0 mm.
*   **Clearance:** 1.0 mm.

Please reply **'confirm'** to apply this change, or **'keep existing'** to keep the current values.
\item User: confirm
\item Assistant: Configuration complete. The following items are confirmed: shell clearance, nested core half length, nested core radius, and shell half length.
\end{itemize}
\textbf{Internal Chain (Natural Language):}
\begin{itemize}
\item Turn 1: action=summarize\_progress, phase=materials, complete=False.
\item Missing fields: materials.volume\_material\_map, source.type, source.particle, source.energy, source.position, source.direction, output.format, output.path, geometry.ask.shell.inner\_radius, geometry.ask.shell.thicknesses, geometry.ask.shell.half\_length, geometry.ask.shell.child\_radius, geometry.ask.shell.child\_half\_length, geometry.ask.shell.clearance.
\item LLM used=True; fallback=none; backend=llm\_semantic\_frame+runtime\_semantic.
\item Turn 2: action=confirm\_overwrite, phase=source, complete=False.
\item Missing fields: none.
\item LLM used=True; fallback=none; backend=llm\_slot\_frame+runtime\_semantic.
\item Turn 3: action=finalize, phase=physics, complete=True.
\item Missing fields: none.
\item LLM used=True; fallback=none; backend=llm\_semantic\_frame+runtime\_semantic+confirmed\_pending\_overwrite.
\end{itemize}
\subsection{Case MTM18 (boolean\_cavity)}
\textbf{Dialogue:}
\begin{itemize}
\item User: 我想做一个挖空的盒体几何，但先只给你意图。
\item Assistant: 好的，我们已确认了您要创建“挖空的盒体”几何结构，并记录了盒体的尺寸。

接下来，为了完成这个布尔减运算的几何体，还需要您提供被挖去的那个形状（Solid B）在 x 和 y 方向上的尺寸。此外，整个几何体的材料也需要您来指定。

请随时告诉我这些信息。
\item User: ???????A ? 40 mm x 40 mm x 40 mm ???B ? 20 mm x 20 mm x 20 mm ???? A ??? B??? ???? gamma 0.3 MeV??? (0,0,-60) mm??? +z????? QBBC??? xml?
\item Assistant: 我注意到，您这次提供的尺寸信息会让几何结构从“挖空盒体”（布尔减运算）变为“合并盒体”（布尔并运算），这会覆盖之前已确认的几何设置。

具体会修改：
- 几何类型：从“挖空盒体”改为“合并盒体”。
- 两个盒体的尺寸也会更新为您新提供的数值。

如果您确认要应用这些修改，请直接回复“确认”；如果您希望保留原来的“挖空盒体”几何，请回复“保持原值”。
\item User: 确认
\item Assistant: 好的，我们已确认了“合并盒体”的几何尺寸、源项、物理列表和输出设置。

接下来，只需要为这个几何体中的各个体积指定材料，就可以生成完整的模拟程序了。
\end{itemize}
\textbf{Internal Chain (Natural Language):}
\begin{itemize}
\item Turn 1: action=summarize\_progress, phase=geometry, complete=False.
\item Missing fields: geometry.params.bool\_a\_x, geometry.params.bool\_a\_y, geometry.params.bool\_a\_z, geometry.params.bool\_b\_x, geometry.params.bool\_b\_y, geometry.params.bool\_b\_z, materials.volume\_material\_map, source.type, source.particle, source.energy, source.position, source.direction, physics.physics\_list, output.format, output.path, geometry.ask.boolean.solid\_a\_size, geometry.ask.boolean.solid\_b\_size.
\item LLM used=True; fallback=none; backend=llm\_slot\_frame+runtime\_semantic.
\item Turn 2: action=confirm\_overwrite, phase=materials, complete=False.
\item Missing fields: materials.volume\_material\_map.
\item LLM used=True; fallback=none; backend=llm\_slot\_frame+runtime\_semantic.
\item Turn 3: action=summarize\_progress, phase=materials, complete=False.
\item Missing fields: materials.volume\_material\_map.
\item LLM used=True; fallback=none; backend=llm\_semantic\_frame+runtime\_semantic+confirmed\_pending\_overwrite.
\end{itemize}
\subsection{Case MTM19 (ring\_complete)}
\textbf{Dialogue:}
\begin{itemize}
\item User: Create a complete ring setup: 16 modules, each 6 mm x 6 mm x 2 mm, radius 32 mm, clearance 0.5 mm, material silicon; point gamma source 0.511 MeV at (0,0,0) mm; physics QBBC; output json.
\item Assistant: Here's your setup summary:

**Completed this turn:**
- **Geometry:** Structure, box X/Y sizes.
- **Materials:** Selected materials and their assignment to volumes.
- **Source:** Type, particle, and energy.
- **Physics:** Physics list.
- **Output:** Format and file path.

**Still needed:**
- **Source:** Direction.
\end{itemize}
\textbf{Internal Chain (Natural Language):}
\begin{itemize}
\item Turn 1: action=summarize\_progress, phase=materials, complete=False.
\item Missing fields: source.direction.
\item LLM used=True; fallback=none; backend=llm\_slot\_frame+runtime\_semantic.
\end{itemize}
\subsection{Case MTM20 (grid\_complete)}
\textbf{Dialogue:}
\begin{itemize}
\item User: ???????????3 x 3 ????????? 12 mm x 12 mm x 3 mm?pitch\_x 15 mm?pitch\_y 15 mm??? 1 mm??? ?????? gamma 0.662 MeV??? (0,0,-90) mm??? +z????? QBBC??? root?
\item Assistant: 目前我们已经完成了几何结构、源设置、物理列表和输出配置的同步。接下来，只需要再补充一下体积与材料的绑定信息，就可以继续推进了。
\end{itemize}
\textbf{Internal Chain (Natural Language):}
\begin{itemize}
\item Turn 1: action=summarize\_progress, phase=materials, complete=False.
\item Missing fields: materials.volume\_material\_map.
\item LLM used=True; fallback=none; backend=llm\_slot\_frame+runtime\_semantic.
\end{itemize}
\subsection{Case MTM21 (layered\_detector)}
\textbf{Dialogue:}
\begin{itemize}
\item User: I want a stacked detector slab, but the individual layer thicknesses are not fixed yet.
\item Assistant: You've set up the basic geometry structure for a stacked detector slab. Here's what we have and what's still needed:

**Completed:**
*   **Geometry:** The overall structure is configured as a stacked slab.

**Still Needed:**
*   **Geometry:** The footprint (size) and thicknesses for the individual stack layers.
*   **Materials:** Which material to assign to your detector volume.
*   **Source:** The type of source, the particle to simulate, and its energy.
*   **Physics:** The physics list to use.
*   **Output:** The format and file path for saving your results.
\item User: Use stack footprint 25 mm x 25 mm, thicknesses 1 mm, 2 mm, 1 mm, layer clearance 0.2 mm, outer box 30 mm x 30 mm x 12 mm, container clearance 0.5 mm, material silicon; beam electron 5 MeV from (0,0,-50) mm along +z; physics FTFP\_BERT; output json.
\item Assistant: This would overwrite a geometry value you've already confirmed.

**Pending Change: Geometry**
*   **Stack Footprint:** 16.33 mm x 10.28 mm → 25.0 mm x 25.0 mm
*   **Stack Layer Thicknesses:** [6.68, 1.19, 5.59] mm → [1.0, 2.0, 1.0] mm
*   **Layer Clearance:** 0.73 mm → 0.2 mm
*   **Container Box Size:** 18.07 x 45.52 x 16.66 mm → 30.0 x 30.0 x 12.0 mm
*   **Container Clearance:** 0.87 mm → 0.5 mm

Please reply **'confirm'** to apply these new values, or **'keep existing'** to keep the current geometry.
\item User: confirm
\item Assistant: Your configuration is complete. The following geometry items have been confirmed: stack container clearance, stack container size, stack layer clearance, and stack layer thicknesses.
\end{itemize}
\textbf{Internal Chain (Natural Language):}
\begin{itemize}
\item Turn 1: action=summarize\_progress, phase=materials, complete=False.
\item Missing fields: materials.volume\_material\_map, source.type, source.particle, source.energy, source.position, source.direction, physics.physics\_list, output.format, output.path, geometry.ask.stack.footprint, geometry.ask.stack.thicknesses, geometry.ask.stack.layer\_clearance, geometry.ask.stack.parent\_size, geometry.ask.stack.nest\_clearance.
\item LLM used=True; fallback=none; backend=llm\_semantic\_frame+runtime\_semantic.
\item Turn 2: action=confirm\_overwrite, phase=source, complete=False.
\item Missing fields: none.
\item LLM used=True; fallback=none; backend=llm\_slot\_frame+runtime\_semantic.
\item Turn 3: action=finalize, phase=physics, complete=True.
\item Missing fields: none.
\item LLM used=True; fallback=none; backend=llm\_semantic\_frame+runtime\_semantic+confirmed\_pending\_overwrite.
\end{itemize}
\subsection{Case MTM22 (embedded\_sensor)}
\textbf{Dialogue:}
\begin{itemize}
\item User: 请创建完整的嵌套配置：外盒 80 mm x 80 mm x 80 mm，内嵌铅圆柱半径 15 mm、半长 25 mm，嵌套间隙 1 mm；材料 ?；点源 gamma 1 MeV，位置 (0,0,-100) mm，方向 +z；物理列表 Shielding；输出 csv。
\item Assistant: 配置已完成！当前已确认的设置包括：体积与材料绑定、输出路径、嵌套间隙（clearance）以及圆柱半径。
\end{itemize}
\textbf{Internal Chain (Natural Language):}
\begin{itemize}
\item Turn 1: action=finalize, phase=materials, complete=True.
\item Missing fields: none.
\item LLM used=True; fallback=none; backend=llm\_slot\_frame+runtime\_semantic.
\end{itemize}
\subsection{Case MTM23 (layered\_shield\_complete)}
\textbf{Dialogue:}
\begin{itemize}
\item User: Create a complete shell configuration: inner radius 12 mm, thicknesses 4 mm and 6 mm, half length 35 mm, nested core radius 8 mm, nested core half length 25 mm, clearance 1 mm, material iron; isotropic gamma source 0.8 MeV at (0,0,0) mm; physics QBBC; output hdf5.
\item Assistant: Your configuration is complete. The following items have been confirmed: box dimensions (X, Y, Z) and the volume-to-material binding.
\end{itemize}
\textbf{Internal Chain (Natural Language):}
\begin{itemize}
\item Turn 1: action=finalize, phase=materials, complete=True.
\item Missing fields: none.
\item LLM used=True; fallback=none; backend=llm\_slot\_frame+runtime\_semantic.
\end{itemize}
\subsection{Case MTM24 (boolean\_complete)}
\textbf{Dialogue:}
\begin{itemize}
\item User: ?????????????????A ? 30 mm x 30 mm x 20 mm ???B ? 10 mm x 10 mm x 40 mm ???? A ? union??? ???? gamma 0.5 MeV??? (0,0,-70) mm??? +z????? QBBC??? json?
\item Assistant: 目前我们已经完成了几何、源、物理和输出的基本设置，只剩下材料的“体积与材料绑定”需要补充。请告诉我每个几何体对应的材料，这样配置就完整了。
\end{itemize}
\textbf{Internal Chain (Natural Language):}
\begin{itemize}
\item Turn 1: action=summarize\_progress, phase=materials, complete=False.
\item Missing fields: materials.volume\_material\_map.
\item LLM used=True; fallback=none; backend=llm\_slot\_frame+runtime\_semantic.
\end{itemize}
\end{document}